\section{Principales aportaciones}

Este trabajo presenta un conjunto de aportaciones técnicas que, consideradas en su conjunto, constituyen una plataforma de diagnóstico vehicular OBD-II completa, de código abierto y bajo coste, diferenciada de las soluciones comerciales existentes en varios aspectos relevantes.

\textbf{Implementación directa de la pila de protocolos.} A diferencia de los adaptadores comerciales de tipo ELM327, que abstraen el protocolo OBD-II mediante comandos AT propietarios, el sistema desarrollado implementa directamente las capas ISO-TP y OBD-II sobre hardware de propósito general. Esto proporciona control total sobre los tiempos de comunicación, la gestión de errores y la extensibilidad del sistema, sin depender de firmware cerrado de terceros.

\textbf{Banco de pruebas reproducible basado en Arduino.} El simulador de ECU desarrollado sobre Arduino MKR con shield CAN constituye una herramienta de validación independiente del vehículo real. El simulador implementa un modelo físico del motor, responde a los cuatro modos OBD-II estándar y permite reproducir condiciones adversas del bus CAN de forma controlada y repetible: pérdida de tramas, latencia variable, respuestas negativas esporádicas e inyección aleatoria de códigos de avería mediante interrupción hardware. Esta capacidad de simulación de ruido es una aportación que distingue este trabajo de implementaciones similares disponibles en la literatura.

\textbf{Arquitectura software modular con separación de capas.} La herramienta de diagnóstico implementada en Python sobre Raspberry Pi sigue un diseño basado en interfaces y capas desacopladas mediante inyección de dependencias. Esta arquitectura facilita la sustitución de cualquier componente (transporte, decodificador, logger) sin modificar el resto del sistema, y permite el uso de un transporte  (\textit{mock}) para el desarrollo y las pruebas sin necesidad de hardware CAN.

\textbf{Acceso inalámbrico mediante BLE GATT con logging persistente.} La integración de un servidor BLE GATT sobre Raspberry Pi, combinada con el registro estructurado de datos en SQLite, permite no solo el acceso en tiempo real a los parámetros del vehículo, sino también la consulta posterior de sesiones históricas completas, incluyendo las tramas CAN en hexadecimal correspondientes a cada comando ejecutado.

\textbf{Aplicación móvil multiplataforma.} La aplicación React Native desarrollada opera sobre Android e iOS con una única base de código, implementa el protocolo JSON sobre NUS para la comunicación con la Raspberry Pi y ofrece una experiencia de usuario orientada tanto a la monitorización en tiempo real como al análisis de datos históricos.

\section{Competencias específicas desarrolladas}

La realización de este Trabajo de Fin de Título ha requerido la aplicación y el desarrollo de competencias específicas del Grado de Ingeniería Informática, plan 41, de la Universidad de Las Palmas de Gran Canaria. A continuación se describe la relación de las competencias más directamente implicadas con el trabajo realizado.

\textbf{CI1 — Capacidad para diseñar, desarrollar, evaluar y asegurar la accesibilidad, ergonomía, usabilidad y seguridad de los sistemas, servicios y aplicaciones informáticas, así como de la información que gestionan.} Esta competencia se ha desarrollado en el diseño de la aplicación móvil, prestando atención a la usabilidad de la interfaz de diagnóstico, y en la implementación de mecanismos de seguridad básicos en la comunicación BLE (autenticación por token, timeout de inactividad).

\textbf{CI5 — Capacidad para especificar, diseñar, implementar y evaluar interfaces persona-computador que garanticen la accesibilidad y usabilidad de los sistemas, servicios y aplicaciones informáticas.} Esta competencia se ha aplicado en el diseño e implementación de la aplicación móvil React Native, con sus cinco pantallas funcionales y el sistema de navegación por pestañas.

\textbf{CI8 — Capacidad para diseñar y evaluar sistemas operativos y servidores, y aplicaciones y sistemas basados en computación distribuida.} El sistema desarrollado es por naturaleza un sistema distribuido: Arduino, Raspberry Pi y dispositivo móvil son tres nodos que se comunican mediante protocolos bien definidos. La Raspberry Pi actúa simultáneamente como cliente CAN y como servidor BLE, gestionando concurrencia mediante hilos y comunicación asíncrona.

\textbf{CI14 — Capacidad para diseñar sistemas de tiempo real y sistemas empotrados.} Esta es la competencia más directamente relacionada con el trabajo. El simulador Arduino opera en tiempo real con interrupciones hardware, temporización precisa de las tramas -TP y gestión concurrente de múltiples tareas (broadcast CAN, respuesta a peticiones OBD-II, generación de ruido). La Raspberry Pi implementa monitorización continua con periodos de muestreo de 500 ms.

\textbf{CI16 — Capacidad para diseñar y desarrollar sistemas de comunicaciones.} El núcleo técnico del trabajo es precisamente la implementación de una pila de protocolos de comunicaciones: CAN (ISO 11898), ISO-TP (ISO 15765-2), OBD-II (SAE J1979) y BLE GATT.

\textbf{CI17 — Capacidad para diseñar y desarrollar sistemas de seguridad informática.} Esta competencia se ha aplicado de forma parcial en el diseño del mecanismo de autenticación BLE y en la gestión segura de las respuestas negativas UDS, que permiten al sistema rechazar peticiones no autorizadas o fuera de contexto.

\textbf{CI18 — Capacidad para diseñar e implementar sistemas de tiempo real y sistemas empotrados.} En complemento con CI14, esta competencia abarca el diseño del flujo de control del simulador Arduino, la implementación del servidor BLE con asyncio en Python y la coordinación entre el hilo de monitorización continua y el hilo de atención a comandos BLE.


\section{Relación con los Objetivos de Desarrollo Sostenible}

Los Objetivos de Desarrollo Sostenible (ODS) de la Agenda 2030 de Naciones Unidas establecen un marco global de referencia para orientar el desarrollo tecnológico hacia un impacto positivo en la sociedad y el medio ambiente. La tabla~\ref{tab:ODS} recoge el grado de relación de este trabajo con cada uno de los 17 ODS.

\begin{table}[!htbp]
\caption{Grado de relación del TFT con los Objetivos de Desarrollo Sostenible de la Agenda 2030.}
\label{tab:ODS}
\centering
\renewcommand{\arraystretch}{1.2}
\begin{tabular}{|l|c|c|c|c|}
\cline{1-5}
 & \multicolumn{4}{c|}{Grado de relación} \\ \cline{2-5}
\textbf{ODS} & \textbf{0} & \textbf{1} & \textbf{2} & \textbf{3} \\
 & No procede & Bajo & Medio & Alto \\ \cline{1-5}
1\quad Fin de la Pobreza              & X & & & \\ \cline{1-5}
2\quad Hambre Cero                    & X & & & \\ \cline{1-5}
3\quad Salud y Bienestar              & & X & & \\ \cline{1-5}
4\quad Educación de Calidad           & & & X & \\ \cline{1-5}
5\quad Igualdad de Género             & X & & & \\ \cline{1-5}
6\quad Agua Limpia y Saneamiento      & X & & & \\ \cline{1-5}
7\quad Energía Asequible y No Contaminante & & & X & \\ \cline{1-5}
8\quad Trabajo Decente y Crecimiento Económico & & X & & \\ \cline{1-5}
9\quad Industria, Innovación e Infraestructuras & & & & X \\ \cline{1-5}
10\quad Reducción de las Desigualdades & X & & & \\ \cline{1-5}
11\quad Ciudades y Comunidades Sostenibles & & X & & \\ \cline{1-5}
12\quad Producción y Consumo Responsables & & & X & \\ \cline{1-5}
13\quad Acción por el Clima           & & & X & \\ \cline{1-5}
14\quad Vida Submarina                & X & & & \\ \cline{1-5}
15\quad Vida de Ecosistemas Terrestres & X & & & \\ \cline{1-5}
16\quad Paz, Justicia e Instituciones Sólidas & X & & & \\ \cline{1-5}
17\quad Alianzas para Lograr los Objetivos & X & & & \\ \cline{1-5}
\end{tabular}
\end{table}


A continuación se justifica el grado de relación asignado a los ODS con mayor vinculación con el trabajo.

\textbf{ODS 9 — Industria, Innovación e Infraestructuras (Alto).} El trabajo desarrolla tecnología e innovación aplicada al sector del mantenimiento automovilístico, un ámbito industrial con un impacto económico significativo. La implementación directa de protocolos de diagnóstico sobre hardware de bajo coste abre posibilidades de innovación en entornos donde las herramientas profesionales resultan inaccesibles económicamente.

\textbf{ODS 12 — Producción y Consumo Responsables (Medio).} El acceso asequible al diagnóstico vehicular fomenta el denominado \textit{derecho a reparar}: un diagnóstico preciso y oportuno permite al propietario del vehículo tomar decisiones informadas sobre el mantenimiento, alargando la vida útil del vehículo y reduciendo la generación de residuos tecnológicos derivados de la sustitución prematura de componentes.

\textbf{ODS 7 y 13 — Energía Asequible y Acción por el Clima (Medio).} El estándar OBD-II incluye parámetros directamente relacionados con la eficiencia de la combustión y las emisiones de gases contaminantes, como el estado del sistema de combustible, los recortes de mezcla (\textit{fuel trim}) o el estado del catalizador. El acceso a esta información facilita la detección temprana de averías que incrementan las emisiones, contribuyendo indirectamente a la reducción de la huella de carbono del transporte.

\textbf{ODS 4 — Educación de Calidad (Medio).} El carácter abierto y documentado del sistema desarrollado, junto con el banco de pruebas basado en Arduino, lo convierte en una herramienta didáctica de alto valor para la enseñanza de protocolos de comunicación industrial, sistemas empotrados y desarrollo de software en capas, en el contexto de titulaciones de ingeniería informática y electrónica.
